\documentclass[a4paper,11pt]{extarticle}


\usepackage[T1]{fontenc}
\usepackage[utf8]{inputenc}
\usepackage{lmodern}

\usepackage[]{parskip}
\usepackage{hyperref, wrapfig, float, amsmath, adjustbox, array, pdflscape, soul, color, extsizes}
\usepackage{booktabs}
\usepackage{lscape}
\usepackage{setspace}

\renewcommand{\arraystretch}{2}
\renewcommand{\contentsname}{Spis treści}
\renewcommand{\refname}{Bibliografia}
\renewcommand{\figurename}{Rys.}
\renewcommand{\tablename}{Tab.}

\usepackage{geometry}
\geometry{
  a4paper,
  margin=30mm,
  bottom=30mm
}

\usepackage{graphicx}
\usepackage[font=footnotesize,labelfont=bf,justification=centering]{caption}
\usepackage{subcaption}
\graphicspath{ {./images/}{../images/} }

\newcommand{\mathcolorbox}[1]{
  \begin{center}
    \colorbox{yellow}{\(\displaystyle #1\)}
  \end{center}
}
\newcommand{\pvec}[1]{
  \vec{#1}\mkern2mu\vphantom{#1}
}
\newcommand{\angstrom}{\mbox{\normalfont\AA}}

\begin{document}
  \begin{titlepage}
    \begin{center}
      \begin{figure}[H]
        \centering
        
\includegraphics[scale=.125]{logopp.png}
      \end{figure}
      \vspace*{.05\textheight}
      \huge\textbf{Politechnika Poznańska}\\
      \large\textbf{Wydział Inżynierii Materiałowej i Fizyki Technicznej}\\
      \large\textbf{Fizyka Techniczna}\\
      \vspace*{.2\textheight}
      \huge\textbf{Detektory gazowe\\promieniowania jonizującego}\\
      \vspace*{.025\textheight}
      \large\textbf{Jakub Bąbelek\\jbbabelek@gmail.com\\(nr.\ albumu: 136454)}
    \end{center}
  \end{titlepage}

  \doublespacing

  %\tableofcontents
  
  \clearpage

  \section{Wstęp}
  Detektory gazowe promieniowania jonizującego to przyrządy wykorzystywane w
  fizyce cząstek elementarnych do wykrywania obecności promieniowania jonizującego
  oraz w ochronie radiologicznej do jego pomiaru. Wykorzystują one efekt jonizacji
  w sensorze wypełnionym gazem. Jeżeli cząstka posiada wystarczającą energię do
  zjonizowania atomu lub molekuły gazu, powstałe elektrony i jony powodują
  mierzalny przepływ napięcia.

  \subsection{Typy detektorów gazowych}
  \begin{itemize}
    \item Komora jonizacyjna
    \item Licznik proporcjonalny
    \item Licznik Geigera
  \end{itemize}

  \subsection{Licznik Geigera}
  Instrument pomiarowy wykorzystywany do detekcji promieniowania jonizującego.
  Powszechnie wykorzystywany w dozymetrii, ochronie radiologicznej, fizyce
  eksperymentalnej czy energetyce jądrowej. Wykrywa promienie \( \alpha \),
  \( \beta \) oraz promienie gamma. Zbudowany jest on z dwóch elektrod w szklanym
  cylindrze wypełnionym mieszaniną gazów (najczęściej stosuje się mieszankę
  argonu i par alkoholu). Trafiająca do wnętrza licznika cząstka wywołuje jonizację
  atomów gazu a w efekcie wyładowanie lawinowe i powstaniem impulsu napięcia
  kierowanym do układu pomiarowego. Zlicza on impulsy i przekształca je w sygnały
  dźwiękowe, świetlne lub na wskazania wskaźnika wychyłowego lub cyfrowego.

  Pomiary w doświadczeniu wykonane zostały przy pomocy licznika Geigera BOH-45.
  Napięcie jego pracy obliczamy ze wzoru:
  
  \begin{equation}
    U_{pracy} = U_{progu} + 100[V]
    \label{eq:napiecie_pracy}
  \end{equation}

  \clearpage
  \section{Dane i obliczenia}

  \begin{table}[H]
    \centering
    \begin{adjustbox}{max width = \textwidth}
    \begin{tabular}{cccc}
      \toprule
      \textbf{Napięcie \( [V] \)} & \textbf{Pomiar 1 \( [\frac{imp}{100s}] \)} & \textbf{Pomiar 2 \( [\frac{imp}{100s}] \)} & \textbf{Pomiar 3 \( [\frac{imp}{100s}] \)} \\
      \midrule
      400                       & 0                 & 0                 & 0                 \\
      420                       & 0                 & 0                 & 0                 \\
      440                       & 13624             & 13506             & 13835             \\
      460                       & 14463             & 14385             & 14456             \\
      480                       & 14650             & 14892             & 14734             \\
      500                       & 14962             & 15071             & 14893             \\
      520                       & 15050             & 15129             & 15284             \\
      540                       & 15209             & 15416             & 15134             \\
      560                       & 15511             & 15519             & 15478             \\
      580                       & 15573             & 15590             & 15401             \\
      600                       & 15801             & 15824             & 15364             \\
      620                       & 15961             & 15706             & 15934             \\
      \bottomrule
    \end{tabular}
    \end{adjustbox}
  \end{table}

  \begin{table}[H]
    \centering
    \begin{adjustbox}{max width = \textwidth}
    \begin{tabular}{ccc}
      \toprule
      \textbf{Napięcie \( [V] \)} & \textbf{Średnia ilość impulsów} & \textbf{Odchylenie standardowe} \\
      \midrule
      400 & 0     & 0   \\
      420 & 0     & 0   \\
      440 & 13655 & 220 \\
      460 & 14435 & 57  \\
      480 & 14759 & 162 \\
      500 & 14975 & 119 \\
      520 & 15154 & 157 \\
      540 & 15253 & 193 \\
      560 & 15503 & 29  \\
      580 & 15521 & 138 \\
      600 & 15663 & 342 \\
      620 & 15867 & 185 \\
      \bottomrule
    \end{tabular}
    \end{adjustbox}
  \end{table}

  Wykorzystująć wzór (\ref{eq:napiecie_pracy}) obliczam \( U_{pracy} \) które wynosi
  \( 540[V] \)

  \begin{figure}[H]
    \centering
    \includegraphics[width=15cm]{wykres.png}
    \caption{Zależnosc średniej ilości impulsów od napięcia}
    \label{fig:wykres}
  \end{figure}

  \section{Wnioski}
  Pomiary pozwoliły ustalić napięcie pracy licznika które wynosi \( 540[V] \).
  Widoczne plateau jest charakterystyczne dla licznika Geigera, formuje się po
  osiągnięciu napięcia progowego.
  Brak pomiarów dla wyższych wartości napięcia nie możemy określić jakości licznika.

\end{document}